% GENERATED FILE. Edit canonical lessons, then run npm run generate:books.
% canonical-source-hash: fnv1a64:1df7bda67ff7d7eb
% canonical-lessons: GE-C526-staub, GE-C526-wischen, GE-C526-fegen, GE-C526-staubsaugen, GE-C526-lappen

\chapter{Wischen und fegen}
\label{ch:ge-wischen-und-fegen}
\hlchaptermodality{526}

\begin{chapteropening}
\textbf{By the end of this chapter:} \emph{I can talk about cleaning the house.}

Complete the last lesson of chapter 526: I can talk about cleaning the house.
\end{chapteropening}

\section[der Staub]{der Staub — dust}

\label{lesson:GE-C526-staub}



\begin{quote}
Before the new one: say the German for to play video games, then the German for headphones.
\end{quote}

\subsection*{You'll want to know: der Staub}
\textbf{der Staub} — \textquotedblleft{}dust\textquotedblright{}.

\begin{quote}
Auf dem alten Schrank liegt viel Staub.
\end{quote}

\begin{quote}
\emph{There's a lot of dust on the old cupboard.}
\end{quote}

No plural.

\begin{cousinweb}
\emph{Staub} is inherited from Old High German \emph{stoup} \textquotedblleft{}dust\textquotedblright{} (EtymDB).
\end{cousinweb}

\begin{grammarlens}[title={What you've built}]
One more word for telling, explaining and describing at B1.
\end{grammarlens}

\subsection*{Your turn}
\begin{itemize}
\raggedright
  \item \emph{Say it:} \emph{der Staub}
  \item \emph{Say it:} \emph{der Staub}, once more
  \item \emph{Say it:} a sentence of your own with \emph{der Staub}
  \item \emph{Recall:} say the German for to play video games, then the German for headphones, then say \emph{der Staub} again
\end{itemize}

\begin{tcolorbox}[breakable,colback=teal!4,colframe=teal!35!black,title={Before you move on}]
What does der Staub mean? (\textquotedblleft{}Dust\textquotedblright{}.) Say it once more.
\end{tcolorbox}
\section[wischen]{wischen — to wipe; to mop}

\label{lesson:GE-C526-wischen}



\begin{quote}
Before the new one: say the German for headphones, then the German for dust.
\end{quote}

\subsection*{You'll want to know: wischen}
\textbf{wischen} — \textquotedblleft{}to wipe; to mop\textquotedblright{}.

\begin{quote}
Am Samstag wische ich den Boden in der Küche.
\end{quote}

\begin{quote}
\emph{On Saturday I mop the kitchen floor.}
\end{quote}

\emph{Staub wischen} — to dust.

\begin{cousinweb}
\emph{wischen} is inherited from Middle High German \emph{wischen} and Old High German \emph{wisken}, from Proto-Germanic \emph{wiskijaną} (EtymDB).
\end{cousinweb}

\begin{grammarlens}[title={What you've built}]
One more word for telling, explaining and describing at B1.
\end{grammarlens}

\subsection*{Your turn}
\begin{itemize}
\raggedright
  \item \emph{Say it:} \emph{wischen}
  \item \emph{Say it:} \emph{wischen}, once more
  \item \emph{Say it:} a sentence of your own with \emph{wischen}
  \item \emph{Recall:} say the German for headphones, then the German for dust, then say \emph{wischen} again
\end{itemize}

\begin{tcolorbox}[breakable,colback=teal!4,colframe=teal!35!black,title={Before you move on}]
What does wischen mean? (\textquotedblleft{}To wipe; to mop\textquotedblright{}.) Say it once more.
\end{tcolorbox}
\section[fegen]{fegen — to sweep}

\label{lesson:GE-C526-fegen}



\begin{quote}
Before the new one: say the German for dust, then the German for to wipe.
\end{quote}

\subsection*{You'll want to know: fegen}
\textbf{fegen} — \textquotedblleft{}to sweep\textquotedblright{}.

\begin{quote}
Mein Vater fegt jeden Morgen die Terrasse.
\end{quote}

\begin{quote}
\emph{My father sweeps the terrace every morning.}
\end{quote}

In southern Germany and Austria also \emph{kehren}.

\begin{cousinweb}
\emph{fegen} goes back to Middle High German \emph{vegen} \textquotedblleft{}to clean, to polish\textquotedblright{}; EtymDB lists the word, glossed \textquotedblleft{}to cleanse, to sweep\textquotedblright{}, without a source.
\end{cousinweb}

\begin{grammarlens}[title={What you've built}]
One more word for telling, explaining and describing at B1.
\end{grammarlens}

\subsection*{Your turn}
\begin{itemize}
\raggedright
  \item \emph{Say it:} \emph{fegen}
  \item \emph{Say it:} \emph{fegen}, once more
  \item \emph{Say it:} a sentence of your own with \emph{fegen}
  \item \emph{Recall:} say the German for dust, then the German for to wipe, then say \emph{fegen} again
\end{itemize}

\begin{tcolorbox}[breakable,colback=teal!4,colframe=teal!35!black,title={Before you move on}]
What does fegen mean? (\textquotedblleft{}To sweep\textquotedblright{}.) Say it once more.
\end{tcolorbox}
\section[staubsaugen]{staubsaugen — to vacuum}

\label{lesson:GE-C526-staubsaugen}



\begin{quote}
Before the new one: say the German for to wipe, then the German for to sweep.
\end{quote}

\subsection*{You'll want to know: staubsaugen}
\textbf{staubsaugen} — \textquotedblleft{}to vacuum\textquotedblright{}.

\begin{quote}
Jeden Freitag staubsaugt mein Mann die ganze Wohnung.
\end{quote}

\begin{quote}
\emph{Every Friday my husband vacuums the whole flat.}
\end{quote}

Both forms are correct: \emph{er staubsaugt} and \emph{er saugt Staub}.

\begin{cousinweb}
\emph{staubsaugen} is built on \emph{Staub} \textquotedblleft{}dust\textquotedblright{} and \emph{saugen} \textquotedblleft{}to suck\textquotedblright{}, inherited from Old High German \emph{sūgan}, from Proto-Germanic \emph{sūganą} (EtymDB).
\end{cousinweb}

\begin{grammarlens}[title={What you've built}]
One more word for telling, explaining and describing at B1.
\end{grammarlens}

\subsection*{Your turn}
\begin{itemize}
\raggedright
  \item \emph{Say it:} \emph{staubsaugen}
  \item \emph{Say it:} \emph{staubsaugen}, once more
  \item \emph{Say it:} a sentence of your own with \emph{staubsaugen}
  \item \emph{Recall:} say the German for to wipe, then the German for to sweep, then say \emph{staubsaugen} again
\end{itemize}

\begin{tcolorbox}[breakable,colback=teal!4,colframe=teal!35!black,title={Before you move on}]
What does staubsaugen mean? (\textquotedblleft{}To vacuum\textquotedblright{}.) Say it once more.
\end{tcolorbox}
\section[der Lappen]{der Lappen — cloth, rag}

\label{lesson:GE-C526-lappen}



\begin{quote}
Before the new one: say the German for to sweep, then the German for to vacuum.
\end{quote}

\subsection*{You'll want to know: der Lappen}
\textbf{der Lappen} — \textquotedblleft{}cloth, rag\textquotedblright{}.

\begin{quote}
Sie putzt das Fenster mit einem feuchten Lappen.
\end{quote}

\begin{quote}
\emph{She cleans the window with a damp cloth.}
\end{quote}

Plural: \emph{die Lappen}. Informal for a driving licence.

\begin{cousinweb}
\emph{Lappen} is inherited from Middle High German \emph{lappe} and Old High German \emph{lappo}, from Proto-Germanic \emph{lappa-} (EtymDB).
\end{cousinweb}

\begin{grammarlens}[title={What you've built}]
One more word for telling, explaining and describing at B1.
\end{grammarlens}

\subsection*{Your turn}
\begin{itemize}
\raggedright
  \item \emph{Say it:} \emph{der Lappen}
  \item \emph{Say it:} \emph{der Lappen}, once more
  \item \emph{Say it:} a sentence of your own with \emph{der Lappen}
  \item \emph{Recall:} say the German for to sweep, then the German for to vacuum, then say \emph{der Lappen} again
\end{itemize}

\begin{tcolorbox}[breakable,colback=teal!4,colframe=teal!35!black,title={Before you move on}]
What does der Lappen mean? (\textquotedblleft{}Cloth, rag\textquotedblright{}.) Say it once more.
\end{tcolorbox}
